% GENERATED FILE. Edit canonical lessons, then run npm run generate:books.
% canonical-source-hash: fnv1a64:0635cdbb88e537da
% canonical-lessons: KA-C394-rajakiya, KA-C394-rajakarani, KA-C394-paksa, KA-C394-bahiskarisu, KA-C394-nayakatva

\chapter{Politics and Leaders}
\label{ch:ka-politics-and-leaders}
\hlchaptermodality{394}

\begin{chapteropening}
\textbf{By the end of this chapter:} \emph{I can say politics, a politician, a political party, to boycott and leadership.}

Complete the last lesson of chapter 394: I can say politics, a politician, a political party, to boycott and leadership.
\end{chapteropening}

\section[\texorpdfstring{\kn{ರಾಜಕೀಯ} (rājakīya)}{ರಾಜಕೀಯ (rājakīya)}]{\kn{ರಾಜಕೀಯ} (rājakīya) — politics; political}

\label{lesson:KA-C394-rajakiya}



\begin{quote}
Before the new one: say the Kannada for democracy, then the Kannada for to represent.
\end{quote}

\subsection*{You'll want to know: \kn{ರಾಜಕೀಯ}}
\textbf{\kn{ರಾಜಕೀಯ}} — \emph{rājakīya} — \textquotedblleft{}politics; political\textquotedblright{}.

\begin{quote}
\kn{ನನ್ನ} \kn{ಅಜ್ಜನಿಗೆ} \kn{ರಾಜಕೀಯದಲ್ಲಿ} \kn{ಆಸಕ್ತಿ} \kn{ಇಲ್ಲ}.
\end{quote}

\begin{quote}
\emph{nanna ajjanige rājakīyadalli āsakti illa.}
\end{quote}

\begin{quote}
My grandfather has no interest in politics.
\end{quote}

As an adjective: \kn{ರಾಜಕೀಯ} \kn{ಸುದ್ದಿ}, political news.

\begin{cousinweb}
\textbf{\kn{ರಾಜಕೀಯ}} (\emph{rājakīya}) is borrowed from Sanskrit \emph{rājakīya}, \textquotedblleft{}belonging to a king, royal\textquotedblright{} (Monier-Williams); Kannada uses it for politics.
\end{cousinweb}

\begin{grammarlens}[title={What you've built}]
One more word for giving reasons and opinions.
\end{grammarlens}

\subsection*{Your turn}
\begin{itemize}
\raggedright
  \item \emph{Say it:} \emph{rājakīya}
  \item \emph{Say it:} \emph{rājakīya}, once more
  \item \emph{Say it:} \emph{nanna ajjanige rājakīyadalli āsakti illa.} — My grandfather has no interest in politics.
  \item \emph{Recall:} say the Kannada for democracy, then the Kannada for to represent, then say \emph{rājakīya} again
\end{itemize}

\begin{tcolorbox}[breakable,colback=teal!4,colframe=teal!35!black,title={Before you move on}]
What does \kn{ರಾಜಕೀಯ} mean? (\textquotedblleft{}Politics; political\textquotedblright{}.) Say it once more.
\end{tcolorbox}
\section[\texorpdfstring{\kn{ರಾಜಕಾರಣಿ} (rājakāraṇi)}{ರಾಜಕಾರಣಿ (rājakāraṇi)}]{\kn{ರಾಜಕಾರಣಿ} (rājakāraṇi) — a politician}

\label{lesson:KA-C394-rajakarani}



\begin{quote}
Before the new one: say the Kannada for to represent, then the Kannada for politics.
\end{quote}

\subsection*{You'll want to know: \kn{ರಾಜಕಾರಣಿ}}
\textbf{\kn{ರಾಜಕಾರಣಿ}} — \emph{rājakāraṇi} — \textquotedblleft{}a politician\textquotedblright{}.

\begin{quote}
\kn{ಆ} \kn{ರಾಜಕಾರಣಿ} \kn{ಹಳ್ಳಿಗಳಿಗೆ} \kn{ಹೋಗಿ} \kn{ಜನರ} \kn{ಜೊತೆ} \kn{ಮಾತಾಡಿದನು}.
\end{quote}

\begin{quote}
\emph{ā rājakāraṇi haḷḷigaḷige hōgi janara jote mātāḍidanu.}
\end{quote}

\begin{quote}
That politician went to the villages and talked with people.
\end{quote}

Plural: \kn{ರಾಜಕಾರಣಿಗಳು}, \emph{rājakāraṇigaḷu}.

\kn{ರಾಜಕಾರಣ}, politics, + -\kn{ಇ}, the one who does it.

\begin{cousinweb}
\textbf{\kn{ರಾಜಕಾರಣಿ}} (\emph{rājakāraṇi}) is a modern Kannada word built on \emph{rājakāraṇa}, \textquotedblleft{}politics\textquotedblright{}, from Sanskrit \emph{rāja}, \textquotedblleft{}king\textquotedblright{}, and \emph{kāraṇa}, \textquotedblleft{}cause, means\textquotedblright{} (CDIAL 3057).
\end{cousinweb}

\begin{grammarlens}[title={What you've built}]
One more word for giving reasons and opinions.
\end{grammarlens}

\subsection*{Your turn}
\begin{itemize}
\raggedright
  \item \emph{Say it:} \emph{rājakāraṇi}
  \item \emph{Say it:} \emph{rājakāraṇi}, once more
  \item \emph{Say it:} \emph{ā rājakāraṇi haḷḷigaḷige hōgi janara jote mātāḍidanu.} — That politician went to the villages and talked with people.
  \item \emph{Recall:} say the Kannada for to represent, then the Kannada for politics, then say \emph{rājakāraṇi} again
\end{itemize}

\begin{tcolorbox}[breakable,colback=teal!4,colframe=teal!35!black,title={Before you move on}]
What does \kn{ರಾಜಕಾರಣಿ} mean? (\textquotedblleft{}A politician\textquotedblright{}.) Say it once more.
\end{tcolorbox}
\section[\texorpdfstring{\kn{ಪಕ್ಷ} (pakṣa)}{ಪಕ್ಷ (pakṣa)}]{\kn{ಪಕ್ಷ} (pakṣa) — a political party; a side}

\label{lesson:KA-C394-paksa}



\begin{quote}
Before the new one: say the Kannada for politics, then the Kannada for a politician.
\end{quote}

\subsection*{You'll want to know: \kn{ಪಕ್ಷ}}
\textbf{\kn{ಪಕ್ಷ}} — \emph{pakṣa} — \textquotedblleft{}a political party; a side\textquotedblright{}.

\begin{quote}
\kn{ಚುನಾವಣೆಯಲ್ಲಿ} \kn{ನಾಲ್ಕು} \kn{ಪಕ್ಷಗಳು} \kn{ಸ್ಪರ್ಧಿಸಿದವು}.
\end{quote}

\begin{quote}
\emph{cunāvaṇeyalli nālku pakṣagaḷu spardhisidavu.}
\end{quote}

\begin{quote}
Four parties contested the election.
\end{quote}

Plural: \kn{ಪಕ್ಷಗಳು}, \emph{pakṣagaḷu}.

Also a side: \kn{ಯಾರ} \kn{ಪಕ್ಷ}?, whose side?

\begin{cousinweb}
\textbf{\kn{ಪಕ್ಷ}} (\emph{pakṣa}) is borrowed from Sanskrit \emph{pakṣa}, \textquotedblleft{}a wing, a side\textquotedblright{} (CDIAL 7627).
\end{cousinweb}

\begin{grammarlens}[title={What you've built}]
One more word for giving reasons and opinions.
\end{grammarlens}

\subsection*{Your turn}
\begin{itemize}
\raggedright
  \item \emph{Say it:} \emph{pakṣa}
  \item \emph{Say it:} \emph{pakṣa}, once more
  \item \emph{Say it:} \emph{cunāvaṇeyalli nālku pakṣagaḷu spardhisidavu.} — Four parties contested the election.
  \item \emph{Recall:} say the Kannada for politics, then the Kannada for a politician, then say \emph{pakṣa} again
\end{itemize}

\begin{tcolorbox}[breakable,colback=teal!4,colframe=teal!35!black,title={Before you move on}]
What does \kn{ಪಕ್ಷ} mean? (\textquotedblleft{}A political party; a side\textquotedblright{}.) Say it once more.
\end{tcolorbox}
\section[\texorpdfstring{\kn{ಬಹಿಷ್ಕರಿಸು} (bahiṣkarisu)}{ಬಹಿಷ್ಕರಿಸು (bahiṣkarisu)}]{\kn{ಬಹಿಷ್ಕರಿಸು} (bahiṣkarisu) — to boycott}

\label{lesson:KA-C394-bahiskarisu}



\begin{quote}
Before the new one: say the Kannada for a politician, then the Kannada for a political party.
\end{quote}

\subsection*{You'll want to know: \kn{ಬಹಿಷ್ಕರಿಸು}}
\textbf{\kn{ಬಹಿಷ್ಕರಿಸು}} — \emph{bahiṣkarisu} — \textquotedblleft{}to boycott\textquotedblright{}.

\begin{quote}
\kn{ಕೆಲಸಗಾರರು} \kn{ಸಭೆಯನ್ನು} \kn{ಬಹಿಷ್ಕರಿಸಿದರು}.
\end{quote}

\begin{quote}
\emph{kelasagāraru sabheyannu bahiṣkarisidaru.}
\end{quote}

\begin{quote}
The workers boycotted the meeting.
\end{quote}

Past (I …): \kn{ಬಹಿಷ್ಕರಿಸಿದೆ}, \emph{bahiṣkariside}.

\kn{ಬಹಿಷ್ಕಾರ}, a boycott.

\begin{cousinweb}
\textbf{\kn{ಬಹಿಷ್ಕರಿಸು}} (\emph{bahiṣkarisu}) is built on Sanskrit \emph{bahiṣkāra}, \textquotedblleft{}expulsion, removal\textquotedblright{} (Monier-Williams), with the Kannada verb ending -\emph{isu}.
\end{cousinweb}

\begin{grammarlens}[title={What you've built}]
One more word for giving reasons and opinions.
\end{grammarlens}

\subsection*{Your turn}
\begin{itemize}
\raggedright
  \item \emph{Say it:} \emph{bahiṣkarisu}
  \item \emph{Say it:} \emph{bahiṣkarisu}, once more
  \item \emph{Say it:} \emph{kelasagāraru sabheyannu bahiṣkarisidaru.} — The workers boycotted the meeting.
  \item \emph{Recall:} say the Kannada for a politician, then the Kannada for a political party, then say \emph{bahiṣkarisu} again
\end{itemize}

\begin{tcolorbox}[breakable,colback=teal!4,colframe=teal!35!black,title={Before you move on}]
What does \kn{ಬಹಿಷ್ಕರಿಸು} mean? (\textquotedblleft{}To boycott\textquotedblright{}.) Say it once more.
\end{tcolorbox}
\section[\texorpdfstring{\kn{ನಾಯಕತ್ವ} (nāyakatva)}{ನಾಯಕತ್ವ (nāyakatva)}]{\kn{ನಾಯಕತ್ವ} (nāyakatva) — leadership}

\label{lesson:KA-C394-nayakatva}



\begin{quote}
Before the new one: say the Kannada for a political party, then the Kannada for to boycott.
\end{quote}

\subsection*{You'll want to know: \kn{ನಾಯಕತ್ವ}}
\textbf{\kn{ನಾಯಕತ್ವ}} — \emph{nāyakatva} — \textquotedblleft{}leadership\textquotedblright{}.

\begin{quote}
\kn{ಅವಳ} \kn{ನಾಯಕತ್ವದಲ್ಲಿ} \kn{ತಂಡ} \kn{ಗೆದ್ದಿತು}.
\end{quote}

\begin{quote}
\emph{avaḷa nāyakatvadalli taṇḍa gedditu.}
\end{quote}

\begin{quote}
Under her leadership the team won.
\end{quote}

Built on the Sanskrit word for a leader, with -\kn{ತ್ವ}, -ship.

\begin{cousinweb}
\textbf{\kn{ನಾಯಕತ್ವ}} (\emph{nāyakatva}) is built on Sanskrit \emph{nāyaka}, \textquotedblleft{}a guide, a leader\textquotedblright{} (Monier-Williams), with -\emph{tva}, \textquotedblleft{}-ship\textquotedblright{}.
\end{cousinweb}

\begin{grammarlens}[title={What you've built}]
One more word for giving reasons and opinions.
\end{grammarlens}

\subsection*{Your turn}
\begin{itemize}
\raggedright
  \item \emph{Say it:} \emph{nāyakatva}
  \item \emph{Say it:} \emph{nāyakatva}, once more
  \item \emph{Say it:} \emph{avaḷa nāyakatvadalli taṇḍa gedditu.} — Under her leadership the team won.
  \item \emph{Recall:} say the Kannada for a political party, then the Kannada for to boycott, then say \emph{nāyakatva} again
\end{itemize}

\begin{tcolorbox}[breakable,colback=teal!4,colframe=teal!35!black,title={Before you move on}]
What does \kn{ನಾಯಕತ್ವ} mean? (\textquotedblleft{}Leadership\textquotedblright{}.) Say it once more.
\end{tcolorbox}
